\documentclass[11pt]{article}
\usepackage[utf8]{inputenc}
\usepackage[english,russian]{babel}
\usepackage{indentfirst}
\usepackage{secdot}
\usepackage{currfile}
\usepackage{datetime}
\setlength{\parindent}{5ex}
\usepackage[left=2.5cm,right=1.5cm,top=2cm,bottom=2cm]{geometry}
\usepackage{fancyhdr}
\title{Лабораторная работ 3}
\author{Рычков Е.А. Федотов Р.Э.}
\date{November 2021}
\begin{document}
\title{\textbf{Подготовка специалистов в области
высокопроизводительных вычислений на базе
межуниверситетской инновационной
учебно-исследовательской лаборатории InterUniLab\footnote{Работа выполняется при финансовой поддержке Министерства образования и науки РФ в рамках прио-ритетного национального проекта «Образование».}}
}
\author{А.С. Абрамова, Н.А. Шехунова, А.В. Бухановский}
\date{}
\maketitle
\thispagestyle{fancy}
\fancyhf{}
\lfoot{\thepage}
\cfoot{\today}
\rfoot{\currfilename}
\begin{abstract}
    Рассматриваются особенности разработки учебно-методического комплекса «Высокопроизводи\-тельные вычисления» на основе модульного и компетентностного под-ходов, ориентированных на слушателей-магистров по специализации «Разработка программного обеспечения» Санкт-Петер- \-бургского Государственного университета информационных технологий, механики и оптики. Самостоятельная работа в рамках комплекса ориентирована на участие в учебно-исследовательских проектах, выпол-няемых в межуниверситетской учебно-исследовательской лаборатории InterUniLab.
\end{abstract}
\section*{\textbf{Введение}}
Современный этап развития высокопроизводительных вычислительных технологий характеризуется: широким распространением многоядерных компьютерных архитектур, удешевлением и доступностью кластерных систем на основе стандартных комплектую-щих, развитием технологий распределенных вычислений, в том числе, Грид \cite{1}. Это тре-бует модификации и развития соответствующих учебных комплексов. Недостаточное внимание сейчас уделяется системному подходу к параллельному математическому и программному обеспечению, как совокупности математических моделей, методов их реа-лизации, параллельных алгоритмов, технологий программирования, тестирования и вери-фикации параллельных программ, хотя именно такой путь позволяет строить эффектив-ные параллельные алгоритмы и проектировать надежные программные системы на их ос-нове \cite{2}.\par
Так как Россия в 2003 году присоединилась к Болонскому процессу, новые УМК должны разрабатываться в соответствии с требованиями, предъявляемыми Европейским союзом. Это позволит включить российские учебные курсы в европейскую систему обра-зования. Главной проблемой перехода к новой системе образования, вызывающей поле-мику, является переход от квалификационного подхода к компетентностному,  а также модульная структура обучения.\par
Преподавание высокопроизводительных вычислений (High Performance Computing, HPC), как дисциплины из области компьютерных наук, требует серьезной материальной и информационной базы. HPC быстро развивается, что приводит к необходимости постоян-ного обновления учебных материалов, которые должны в общем случае содержать мультидис\-циплинарные сведения (архитектура ЭВМ, теория построения алгоритмов, техноло-гии программирования, коммуникационные технологии и пр.). Требования к педагогиче-скому процессу в области высокопроизводительных вычислений, такие как направлен-ность на конкретный результат образования и гибкость, дают право описать его в виде функциональной системы. Описание педагогического процесса в виде функциональной системы дает возможность эффективно управлять этим сложным процессом \cite{3}.
\newpage
\pagestyle{fancy}
\fancyhf{}
\lfoot{\thepage}
\cfoot{\today}
\rfoot{\currfilename}
\section{Разработка курса.}
Курс «Конструирование и анализ параллельных алгоритмов» разбит на модули: про-ектирование параллельных алгоритмов, проектирование параллельных программ, при-кладные параллельные алгоритмы. По окончании курса студент должен уметь: строить эффективные параллельные алгоритмы, применительно к конкретной вычислительной ар-хитектуре, уметь оценивать и моделировать параллельную производительность алгоритма для определенной вычислительной архитектуры, выбирать алгоритмы для решения по-ставленной задачи, выбирать технологии параллельного программирования для решения поставленной задачи, реализовывать ПО с помощью технологий параллельного програм-мирования, оценивать эффективность работы параллельной программы и формулировать рекомендации по ее модификации.
\par Курс «Технологии распределенный вычислений и систем» разбит на 5 модулей: ос-новные виды распределенных вычислительных архитектур, концепция и модели Грид, проектирование приложений для распределенных вычислительных архитектур, разработ-ка приложений в peer-to-peer-системах, разработка приложений в современных Грид-системах. По окончании курса студенты должны уметь выбирать оптимальный способ ор-ганизации распределенной вычислительной системы, классифицировать распределенные вычислительные архитектуры, оценивать и моделировать параллельную производитель-ность распределенной вычислительной системы, выбирать вычислительную систему для решения поставленной задачи, выбирать программный инструментарий для решения по-ставленной задачи, пользоваться программным инструментарием распределенных вычис-лительных систем, разрабатывать эффективное программное обеспечение для распреде-ленный вычислительных систем.
\par В методическое обеспечение курса входит виртуальная лаборатория, которая форми-руется на базе разрабатываемых лабораторных работ. Лабораторные работы посвящены построению и оптимизации параллельных алгоритмов: метод Монте-Карло вычисления интегралов, решение систем линейных алгебраических уравнений методом Монте-Карло, генетический алгоритм (глобальная оптимизация), численное интегрирование (квадрату-ры), поиск на графах, умножение матриц, метод конечных элементов, параллельное LU-разложение.
\section{Работа студента.}
Самостоятельная работа студентов в рамках данных курсов ориентирована на участие в учебно- \-исследовательских проектах, выполняемых в межуниверситетской инновационной учебно-исследо \-вательской лаборатории InterUniLab. InterUniLab создана совместной инициативой Санкт-Петербург \-ских университетов — Санкт-Петербургского Государственного политехнического университета, СПбГУ ИТМО, Санкт-Петер\-бургского Госу-дарственного университета авиаприборостроения и др. — и Фондом содействия развитию малых предприятий в научно-технической сфере при поддержке глобальных IT-компаний, таких как Intel, Microsoft, Cadence. Одной из ее задач является подготовка квалифициро-ванных кадров в области критических технологий (в том числе, технологии расп\-ределенных вычислений и систем, высокопроизводительные вычисления) путем вовлечения слу-шателей в практическую реализацию мотивационных (курсовых) проектов — индивиду-ального или в составе рабочей группы. Мотивационный проект ориентирован на разработ-ку математического обеспечения высокопроизводительных вычислений в определенной предметной области. Слушателям на выбор будут предложены задачи из области гидро-метеорологии, экологии, биомедицины, физики плазмы, технической диагностики и управления подвижными техническими объектами, основанные на реальных массивах данных.
\par УМК «Высокопроизводительные вычисления» в СПбГУ ИТМО в настоящий момент находится в состоянии разработки. Однако отдельные его элементы уже прошли апроба-цию в рамках летних и зимних школ Intel (2006, 2007 гг.), а также в плановом учебном процессе СПбГУ ИТМО. Ввод УМК в опытную эксплуатацию планируется в осеннем се-местре 2008 г.
\newpage
\pagestyle{fancy}
\fancyhf{}
\lfoot{\thepage}
\cfoot{\today}
\rfoot{\currfilename}
\section{Транслятор}
Транслятор с языка программирования Аспект имеет классическую архитектуру и состоит из лексического анализа, синтаксического анализа контекстного анализа, генератора внутреннего представления и генератора кода. Реализация первых четырех модулей стандартна и далее не рассматривается.
На вход транслятору подается один или несколько файлов, в каждом из которых может быть одна или несколько А-программ, а также специальный конфигурационный файл, в котором для каждого массового А-блока указан требуемый размер А-фрагмента по каждой размерности. На выходе транслятор генерирует фрагментированную программу на языке C++, состоящую из одного главного файла (main.cpp) и нескольких заголовочных файлов – по одному на каждую найденную А-программу.\par
Каждая А-программа представляется в виде класса, наследуемого от виртуального класса \-AProgram. А-программы могут быть вложенными (операция А-блока может реализовываться другой А-программой), и наличие единого предка позволяет унифицировать управление деревом А-программ.\par
Локальные информационные переменные представляются закрытыми (private) переменными класса, а входные/выходные информационные переменные – закрытыми ссылками. Все управляющие переменные также являются закрытыми. Для доступа к информационным переменным автоматически генерируются inline-функции доступа, которые помимо возврата данных могут производить проверку корректности обращения (например, проверку на выход за границы индекса).\par
Для каждой триггер-функции генерируется открытая (public) функция типа bool, а для каждой операции и управляющего оператора – соответствующие открытые процедуры.
\section{Исполнительная подсистема}
Исполнительная подсистема состоит из трех слоев: слоя абстрагирования от ОС, слоя функциональных модулей и слоя А-фрагментов.\par
Слой абстрагирования от ОС включается в себя все подпрограммы, в которых используются API операционной системы (создание/уничтожение/синхронизация потоков, функции определения доступных ресурсов, функции для работы со временем и т.д.). Это позволяет переносить исполнительную подсистему из одной ОС в другую заменой лишь одного, выделенного слоя.\par
Менеджер процессоров управляет распределением работы между процессорами (ядрами). При старте программы создается по одному потоку на каждый процессор (ядро). Каждый поток обращается за очередной порцией работы в общую очередь, которая реализована в виде монитора Хоара. Программа завершается, когда очередь становится пустой.\par
Менеджер памяти управляет распределением памяти в системе. С точки зрения фрагментированной модели вычислений его основной задаче является представление всех данных во фрагментированном виде, т.е. если массовый А-блок разбивается на n А-фрагментов по k экземпляров в каждом, то данные, обрабатываемые этим А-блоком, также хранятся во фрагментированном виде (n фрагментов по k элементов в каждом). Это позволяется при необходимости легко переносить один или несколько А-фрагментов с одного узла на другой (например, для динамической балансировки загрузки).\par
Менеджер коммуникаций отвечает за прием/передачу сообщений между узлами ЭВМ с распределенной памятью, однако в настоящее время этот модуль не реализован.
\newpage
\pagestyle{fancy}
\fancyhf{}
\lfoot{\thepage}
\cfoot{\today}
\rfoot{\currfilename}
\section{Результаты экспериментов}
К настоящему моменту разработка системы программирования Аспект завершена не полностью и эксперименты по решению практических задач не проводились. Поэтому рассмотрим результаты тестирования ядра самой системы Аспект на примере модельной задачи непроцедурного умножения матриц (текст программы на языке Аспект приведен в разделе 3.2).\par
Для тестирования использовалась следующая конфигурация: Athlon 64 X2 3600+ (2*256 L2), 1024 DDR2 RAM, Windows Vista Ultimate. Результаты тестирования для матриц размером 400 на 400 элементов приведены на рис. 4 (отметка Б/А показывает скорость аналогичной программы, разработанной вручную).\par
Из диаграммы видно, что ускорение достигает значения 1.93 (при размере А-фрагмента 50 экземпляров). Это объясняется тем, что хотя каждое ядро имеет отдельный кэш второго уровня, доступ к оперативной памяти производится через общий контроллер. Таким образом, эффективное использование кэш-памяти для многоядерных процессоров имеет первостепенное значение.\par
Скачок при размере А-фрагментов в 400 экземпляров происходит потому, что его размер совпадает с исходным размером матриц, т.е. вся матрица представляет собой один А-фрагмент. В этой ситуации для второго ядра просто нет работы.\par
Следует обратить внимание, что ось А-фрагментов не линейна. Например, при размере 
А-фрагментов в 100 экземпляров всего получается 80 А-фрагментов (4*4*4 + 4*4), а при размере в 50 экземпляров – уже 576 А-фрагментов. Таким образом, при существенном увеличении числа фрагментов скорость счета не возрастает (она даже несколько снижается под влиянием кэш-памяти), что свидетельствует об эффективности работы ядра системы.\par


\begin{thebibliography}{Литература}
\bibitem{1}Defining the Grid — a snapshot of the current view // H. Stockinger, 2006 ( www.gridclub.ru)
\bibitem{2}Гергель В.П., Стронгин Р.Г. Основы параллельных вычислений для многопроцес-сорных вычислительных систем. – Н.Новгород, ННГУ, 2001
\bibitem{3}Анохин П.К. Принципиальные вопросы общей теории функциональных систем. — М, 1973
\end{thebibliography}
Что произойдёт с текстом документа, если установить значение команды textheight  равным 29,7cm?

Для чего предназначен параметр empty команды pagestyle?
Команда убирает номерование страниц.
\end{document}